\documentclass[a4paper]{article}

\usepackage[a4paper, top=8mm, bottom=8mm, left=8mm, right=8mm]{geometry}

\usepackage{polyglossia}
\setdefaultlanguage[babelshorthands=true]{russian}
\setotherlanguage{english}

\usepackage{fontspec}
\setmainfont{FreeSerif}
\newfontfamily{\russianfonttt}[Scale=0.7]{DejaVuSansMono}

\usepackage[tiny, compact]{titlesec}

\usepackage{titling}
\setlength{\droptitle}{-1cm}
\pretitle{\begin{center}\begin{bfseries}\Large}
\posttitle{\par\end{bfseries}\end{center}}
\preauthor{\begin{center}\normalsize}
\postauthor{\par\end{center}\vspace{-1.8cm}}

\usepackage{hyperref}
\usepackage{bookmark}
\usepackage{csquotes}
% Сюда можно дописывать нужные вам пакеты.

\title{Реализация типа \enquote{Множество}, основанная на красно-чёрном дереве, на языке F\#}

\author{Стариков Егор Юрьевич}

% Дата много места занимает, её лучше не указывать.
\date{}

\begin{document}

\maketitle

\begin{flushright}
    Группа: \emph{\enquote{25.Б42-мм}}

    Кафедра: \emph{системного программирования}

    % Степень/звание и должность научника мы лучше вас знаем, их можно не указывать.
    Научный руководитель: \emph{Григорьев Семён Вячеславович}

    % А вот про консультанта укажите официальное место работы, с "ООО" или чем-то таким, и должностью желательно по трудовой, не просто "techlead". Выясните у консультанта.
    
    
    Номер семестра практики: \emph{2}
\end{flushright} 


\section{Введение}

LamaGraph~--- проект, целью которого является разработка многоядерного вычислителя для разреженной линейной алгебры.
QTreeFSharp представляет собой библиотеку, разрабатываемую в рамках LamaGraph, содержащую структуры данных, необходимые проекту. 
Для расширения её возможностей требовалась реализация структуры данных \enquote{Множество}, базовых операций над ней: вставки, удаления и поиска, 
а также теоретико-множественные операции: объединение, пересечение, разность.

\section{Постановка задачи}
Целью работы является реализация структуры данных \enquote{Множество}, основанная на красно-чёрном дереве для библиотеки QTreeFSharp.

Для достижения данной цели были поставлены следующие задачи: 
\begin{enumerate}
    \item реализовать множество и базовые операции над ним;
    \item реализовать теоретико-множественные операции;
    \item протестировать все реализованные операции;
    \item выполнить сравнение реализованного множества и эталонного из стандартной библиотеки F\#.
\end{enumerate}


\section{Описание решения}
Красно-черное дерево~--- классический способ реализации множества: узел хранит значение, 
ссылки на два своих поддерева и окрашен в цвет(красный или чёрный). 
Его инварианты позволяют эффективно балансировать дерево при операциях вставки и удаления. 
В данной работе использована оптимизированная версия Кэмерона Моя~\cite{Moy}, применяющая 
паттерн \texttt{Done}/\texttt{ToDo} при балансировке. Результат
рекурсивного вызова помечается: нужна ли перебалансировка на обратном ходе
рекурсии. Если поддерево после вставки или удаления не изменило чёрную высоту, 
родительский узел не вызывает \texttt{balance}~--- это
устраняет лишние проверки, ускоряет вставку и упрощает удаление.

Для теоретико-множественных операций реализованы три вспомогательные функции,
работающие не поэлементным обходом, а перестройкой структуры с использованием
инвариантов красно-чёрного дерева:
\begin{enumerate}
    \item \textbf{\texttt{join} $t_1\ g\ t_2$}~--- сшивает два дерева через
          разделитель $g$, про который известно, что все ключи $t_1 < g <$
          всех ключей $t_2$. Функция сравнивает чёрные высоты $h_1$ и $h_2$ и
          спускается по более высокому дереву вдоль одной ветви на разницу
          высот, вставляя второе дерево на нужную глубину и локально
          перебалансируя на обратном ходе. Сложность~--- $O(|h_1 - h_2|) =
          O(\log n)$;

    \item \textbf{\texttt{split} $k_x\ t$}~--- делит дерево на два: $lt$ (все
          ключи $< k_x$) и $gt$ (все ключи $> k_x$). Если ключ $k_x$ найден,
          он исключается из результата. Функция спускается от корня к $k_x$
          по одному пути; на каждом шаге отрезанная часть присоединяется к
          результату через \texttt{join}. Сложность~--- $O(\log n)$;

    \item \textbf{\texttt{merge} $t_1\ t_2$}~--- сшивает два дерева без
          известного разделителя: порядок ключей между $t_1$ и $t_2$ не
          предполагается. Функция находит минимальный элемент во втором
          дереве, удаляет его и использует как разделитель в \texttt{join}.
          Сложность~--- $O(\log n)$.
\end{enumerate}

На основе \texttt{join}, \texttt{split} и \texttt{merge} операции
\texttt{union}, \texttt{intersection} и \texttt{difference} строятся
рекурсивно. Идея состоит в том, что одно дерево разрезается по корню другого
функцией \texttt{split}; полученные независимые левая и правая части
обрабатываются рекурсивно, а результаты сшиваются через \texttt{join} (если
ключ входит в результат) или \texttt{merge} (если ключ не входит). Такой подход
даёт \emph{структурное переиспользование поддеревьев} вместо поэлементной
вставки и удаления.


\section{Эксперименты}

Эксперименты проводились на машине с операционной системой Linux Ubuntu 24.04.3 и процессором Intel Core i7-9750H CPU 2.60GHz с 4 логичискими и 4 физичискими ядрами.
Сравнение проводилось между RBSet и Set путем генерации случайных наборов чисел размером 100, 10000, 100000. 
Стандартный Set в среднем демонстрирует превосходство в несколько раз по таким критериям, как время выполнения операций и используемая память, 
по сравнению с RBSet, кроме отдельных случаев.
\begin{enumerate}
    \item При пакетных вставках/удалениях Set опережает RBSet в \textbf{2.5--5.9 раза}.
          Максимальный выигрыш Set наблюдается при средних размерах
          пакета ($N = 100$--$1000$) и малых деревьях ($A = 100$).
          При больших деревьях ($A = 100\,000$) разница во времени на выполнение 
          операций Set и RBSet сокращается до \textbf{1.5–2 раз}, что можно объяснить 
          доминированием стоимости обхода дерева над стоимостью аллокаций.
           Set аллоцирует в \textbf{3.5--4.5 раза}
          меньше памяти на вставки/удаления по сравнению с RBSet, 
          что указывает на более эффективный
          \textit{structural sharing} при перестроении дерева.
    \item При объединении двух множеств Set стабильно превосходит RBSet в \textbf{10 раз} по времени, необходимому для 
          завершения данной операции.
          Set использует в \textbf{4.5 раза} меньше памяти на объединение двух множеств по сравнению с RBSet.
    \item При пересечении и разности множеств были выявлены следующие результаты: 
          при сильной асимметрии множеств: $A = 100$ и $B = 100\,000$ элементов для разности и 
           $A = 100\,000$ и $B = 100$ для пересечения RBSet опередил Set по времени выполнения операции в \textbf{4.69} и \textbf{5.09} раза соответственно.
           Такая разница объясняется более оптимизированным выбором дерева для обхода в RBSet: join выбирает дерево с меньшей высотой, что позволяет не тратить 
           время на спуск, а также на аллокацию промежуточных поддеревьев, которые Set не может оптимально использовать из-за большой разницы в количестве 
           элементов (большинства элементов большего дерева нет в меньшем). Также при таком соотношении размеров множеств при операции разности RBSet аллоцирует 
           в \textbf{9.7 раз} меньше памяти, чем Set. В остальных случаях Set имеет громадное преимущество, так как в большинстве замеров ответом становится пустое 
           множество, и Set затрачивает только 40 байт, чтобы вернуть его, в то время как RBSet аллоцирует до нескольких сотен килобайт для промежуточных деревьев. 
\end{enumerate}  

\section{Заключение}

Реализация \enquote{Множества} на красно-чёрном дереве не имеет практического смысла, если есть возможность использовать встроенные решения языка, 
во всяком случае, в данной версии. При этом были продемонстрированы недостатки стандартной реализации, касающиеся операций пересечения и разности, 
что может быть использовано в последующих проектах, направленных на разработку оптимальной реализации типа \enquote{Множество}.

В ходе работы получены следующие результаты:
\begin{itemize}
    \item реализована структура данных \enquote{Множество}, основанная на красно-чёрном дереве;
    \item реализованы теоретико-множественные операции;
    \item написаны тесты для всех реализованных операций;
    \item проведены сравнительные замеры производительности и их анализ.
\end{itemize}

Пуллреквест: \url{https://github.com/Lamagraph/QTreeFSharp/pull/16}.



\section{Литература}
\begin{thebibliography}{1}
    \bibitem{Moy} 
    Moy C. Faster, Simpler Red-Black Trees // Trends in Functional Programming. --- Springer, Cham, 2023. --- Vol. 14151. --- P. 36--50. --- DOI: 10.1007/978-3-031-38938-2\_3.
\end{thebibliography}
\end{document}